\documentclass[12pt,a4paper]{article}
        \makeatletter
        \def\input@path{{../}}
        \makeatother
        \usepackage{almeria-fichas}

        \FichaMeta{Sesion 14}{Leer e interpretar graficas}{Bloque 3 · Datos, funciones y graficas}{Identificar maximos, minimos y tendencias para explicar que esta pasando en los datos de la ciudad.}{Conceptual}{Comprender, Analizar, Evaluar}

        \begin{document}

        \FichaCabecera

        \begin{materialesbox}
        \begin{itemize}
    \item Ficha impresa
    \item Lapiz
    \item Regla
    \item Calculadora si se necesita
\end{itemize}
        \end{materialesbox}

        \begin{pasosbox}
        \begin{enumerate}
    \item Lee primero el titulo de la grafica.
    \item Despues mira los ejes y las unidades.
    \item Busca maximos, minimos y tramos crecientes o decrecientes.
    \item Escribe una conclusion final con una frase completa.
\end{enumerate}
        \end{pasosbox}

        \begin{recordatoriobox}
        \begin{itemize}
    \item Maximo: valor mas alto.
    \item Minimo: valor mas bajo.
    \item Una tendencia puede ser creciente, decreciente o estable.
\end{itemize}
        \end{recordatoriobox}

        \begin{apoyobox}
        \begin{itemize}
    \item Señala con una flecha el punto mas alto y el mas bajo.
    \item Haz una tabla pequeña con lo que ves en la grafica.
    \item Si comparas dos anos, escribe primero los valores.
\end{itemize}
        \end{apoyobox}

        \BloomSection{comprender}

\BloomExercise{comprender}{1}{Lectura guiada}{A partir de una grafica de temperatura, explica que significa que julio tenga un valor mayor que abril.}{6}

\BloomExercise{comprender}{2}{Tendencia}{Explica la diferencia entre un tramo creciente y un tramo decreciente.}{5}

\BloomSection{analizar}

\BloomExercise{analizar}{1}{Anomalia}{Analiza por que en una grafica de turismo puede aparecer una bajada fuerte en un ano concreto.}{7}

\BloomExercise{analizar}{2}{Relacion con el contexto}{Una grafica de terrazas abiertas sube cuando sube la temperatura. Analiza que conclusion se puede extraer y que no se puede asegurar del todo.}{7}

\BloomSection{evaluar}

\BloomExercise{evaluar}{1}{Valora una propuesta}{Evalua una posible solucion o producto relacionado con \emph{Leer e interpretar graficas}. Decide si es adecuado y justifica tu opinion con criterios matematicos claros.}{8}

\BloomExercise{evaluar}{2}{Elige la mejor opcion}{Entre varias opciones posibles para cumplir este objetivo: \emph{Identificar maximos, minimos y tendencias para explicar que esta pasando en los datos de la ciudad.}, elige la mejor y argumenta por que. Incluye al menos dos criterios de valoracion.}{8}

        \begin{evidenciabox}
        \begin{itemize}
    \item Leer valores concretos en una grafica.
    \item Describir tendencias con vocabulario adecuado.
    \item Interpretar que puede explicar un cambio notable.
\end{itemize}
        \end{evidenciabox}

        \end{document}
